\subsection{ROOTS OF A SPLIT SEMI-SIMPLE LIE ALGEBRA}

In this number, \((\mathfrak{g},\mathfrak{h})\) denotes a split semi-simple Lie algebra.

For any \(\lambda \in \mathfrak{h}^*\) , denote by \(\mathfrak{g}^{\lambda}(\mathfrak{h})\) , or simply by \(\mathfrak{g}^{\lambda}\) , the primary subspace of \(\mathfrak{g}\) relative to \(\lambda\) (cf.~Chap. VII, §1, no. 3). Recall that \(\mathfrak{g}^{0} = \mathfrak{h}\) (Chap. VII, §2, no. 1, Prop. 4), that \(\mathfrak{g}\) is the direct sum of the \(\mathfrak{g}^{\lambda}\) (Chap. VII, §1, no. 3, Prop. 8 and 9), that \(\mathfrak{g}^{\lambda}\) is the set of \(x \in \mathfrak{g}\) such that \([h, x] = \lambda(h)x\) for all \(h \in \mathfrak{h}\) (Chap. VII, §2, no. 4, Cor. 1 of Th. 2), and that the weights of \(\mathfrak{h}\) on \(\mathfrak{g}\) are the linear forms \(\lambda\) on \(\mathfrak{h}\) such that \(\mathfrak{g}^{\lambda} \neq 0\) (Chap. VII, §1, no. 1).

DEFINITION 2. A root of \((\mathfrak{g}, \mathfrak{h})\) is a non-zero weight of h on g.

Denote by \(R(\mathfrak{g},\mathfrak{h})\) , or simply by R, the set of roots of \((\mathfrak{g},\mathfrak{h})\) . We have

\[
\mathfrak {g} = \mathfrak {h} \oplus \bigoplus_ {\alpha \in R} \mathfrak {g} ^ {\alpha}.
\]

PROPOSITION 1. Let \(\alpha, \beta\) be roots of \((\mathfrak{g}, \mathfrak{h})\) and let \(\langle \cdot, \cdot \rangle\) be a non-degenerate invariant symmetric bilinear form on \(\mathfrak{g}\) (for example the Killing form of \(\mathfrak{g}\) ).

\begin{enumerate}
\def\labelenumi{(\roman{enumi})}
\item
  If \(\alpha + \beta \neq 0\) , \(\mathfrak{g}^{\alpha}\) and \(\mathfrak{g}^{\beta}\) are orthogonal. The restriction of \(\langle \cdot, \cdot \rangle\) to \(\mathfrak{g}^{\alpha} \times \mathfrak{g}^{-\alpha}\) is non-degenerate. The restriction of \(\langle \cdot, \cdot \rangle\) to \(\mathfrak{h}\) is non-degenerate.
\item
  Let \(x \in \mathfrak{g}^{\alpha}\) , \(y \in \mathfrak{g}^{-\alpha}\) and \(h \in \mathfrak{h}\) . Then \([x, y] \in \mathfrak{h}\) and
\end{enumerate}

\[
\langle h, [ x, y ] \rangle = \alpha (h) \langle x, y \rangle .
\]

Assertion (i) is a particular case of Prop. 10 (iii) of Chap. VII, §1, no. 3. If \(x \in \mathfrak{g}^{\alpha}\) , \(y \in \mathfrak{g}^{-\alpha}\) and \(h \in \mathfrak{h}\) , we have \([x, y] \in \mathfrak{g}^{\alpha - \alpha} = \mathfrak{h}\) , and

\[
\langle h, [ x, y ] \rangle = \langle [ h, x ], y \rangle = \langle \alpha (h) x, y \rangle = \alpha (h) \langle x, y \rangle .
\]

THEOREM 1. Let \(\alpha\) be a root of \((\mathfrak{g},\mathfrak{h})\) .

\begin{enumerate}
\def\labelenumi{(\roman{enumi})}
\item
  The vector space \(\mathfrak{g}^{\alpha}\) is of dimension 1.
\item
  The vector subspace \(\mathfrak{h}_{\alpha} = [\mathfrak{g}^{\alpha},\mathfrak{g}^{-\alpha}]\) of \(\mathfrak{h}\) is of dimension 1. It contains a unique element \(H_{\alpha}\) such that \(\alpha (H_{\alpha}) = 2\) .
\item
  The vector subspace \(\mathfrak{s}_{\alpha} = \mathfrak{h}_{\alpha} + \mathfrak{g}^{\alpha} + \mathfrak{g}^{-\alpha}\) is a Lie subalgebra of \(\mathfrak{g}\) .
\item
  If \(X_{\alpha}\) is a non-zero element of \(\mathfrak{g}^{\alpha}\) , there exists a unique \(X_{-\alpha} \in \mathfrak{g}^{-\alpha}\) such that \([X_{\alpha}, X_{-\alpha}] = -H_{\alpha}\) . Let \(\varphi\) be the linear map from \(\mathfrak{sl}(2, k)\) to \(\mathfrak{g}\) that takes \(X_{+}\) to \(X_{\alpha}\) , \(X_{-}\) to \(X_{-\alpha}\) , and \(H\) to \(H_{\alpha}\) ; then \(\varphi\) is an isomorphism from the Lie algebra \(\mathfrak{sl}(2, k)\) to the Lie algebra \(\mathfrak{s}_{\alpha}\) .
\end{enumerate}

\begin{enumerate}
\def\labelenumi{\alph{enumi})}
\item
  Let \(h_{\alpha}\) be the unique element of h such that \(\alpha(h) = \langle h_{\alpha}, h \rangle\) for all \(h \in h\) . By Prop. 1, \([x, y] = \langle x, y \rangle h_{\alpha}\) for all \(x \in g^{\alpha}\) , \(y \in g^{-\alpha}\) ; on the other hand \(\langle g^{\alpha}, g^{-\alpha} \rangle \neq 0\) . Hence \(h_{\alpha} = [g^{\alpha}, g^{-\alpha}] = kh_{\alpha}\) .
\item
  Choose \(x \in \mathfrak{g}^{\alpha}\) , \(y \in \mathfrak{g}^{-\alpha}\) such that \(\langle x, y \rangle = 1\) , so \([x, y] = h_{\alpha}\) . Recall that \([h_{\alpha}, x] = \alpha(h_{\alpha})x\) , \([h_{\alpha}, y] = -\alpha(h_{\alpha})y\) . If \(\alpha(h_{\alpha}) = 0\) , it follows that \(kx + ky + kh_{\alpha}\) is a nilpotent subalgebra \(t\) of \(\mathfrak{g}\) ; since \(h_{\alpha} \in [t, t]\) , \(\operatorname{ad}_{\mathfrak{g}} h_{\alpha}\) is nilpotent (Chap. I, §5, no. 3, Th. 1), which is absurd since \(\operatorname{ad}_{\mathfrak{g}} h_{\alpha}\) is non-zero semi-simple. So \(\alpha(h_{\alpha}) \neq 0\) . Hence there exists a unique \(H_{\alpha} \in \mathfrak{h}_{\alpha}\) such that \(\alpha(H_{\alpha}) = 2\) , which proves (ii).
\item
  Choose a non-zero element \(X_{\alpha}\) of \(\mathfrak{g}^{\alpha}\) . There exists \(X_{-\alpha} \in \mathfrak{g}^{-\alpha}\) such that \([X_{\alpha}, X_{-\alpha}] = -H_{\alpha}\) (since \([X_{\alpha}, \mathfrak{g}^{-\alpha}] = \mathfrak{h}_{\alpha}\) by \(b\) ). Then
\end{enumerate}

\[
\begin{array}{c} [ H _ {\alpha}, X _ {\alpha} ] = \alpha (H _ {\alpha}) X _ {\alpha} = 2 X _ {\alpha}, [ H _ {\alpha}, X _ {- \alpha} ] = - \alpha (H _ {\alpha}) X _ {- \alpha} = - 2 X _ {- \alpha}, \\ \relax [ X _ {\alpha}, X _ {- \alpha} ] = - H _ {\alpha}; \end{array}
\]

hence \(kX_{\alpha} + kX_{-\alpha} + kH_{\alpha}\) is a subalgebra of \(\mathfrak{g}\) and the linear map \(\varphi\) from \(\mathfrak{sl}(2,k)\) to \(kX_{\alpha} + kX_{-\alpha} + kH_{\alpha}\) such that \(\varphi(X_{+}) = X_{\alpha}, \varphi(X_{-}) = X_{-\alpha}, \varphi(H) = H_{\alpha}\) is an isomorphism of Lie algebras.

\begin{enumerate}
\def\labelenumi{\alph{enumi})}
\setcounter{enumi}{3}
\tightlist
\item
  Assume that \(\dim g^{\alpha} > 1\) . Let y be a non-zero element of \(g^{-\alpha}\) . There exists a non-zero element \(X_{\alpha}\) of \(g_{\alpha}\) such that \(\langle y, X_{\alpha} \rangle = 0\) . Choose \(X_{-\alpha}\) as in c), and consider the representation \(\rho : u \mapsto \operatorname{ad}_{\mathfrak{g}} \varphi(u)\) from \(\mathfrak{sl}(2, k)\) to g. We have
\end{enumerate}

\[
\begin{array}{c} \rho (H) y = [ \varphi (H), y ] = [ H _ {\alpha}, y ] = - 2 y \\ \rho (X _ {+}) y = [ \varphi (X _ {+}), y ] = [ X _ {\alpha}, y ] = \langle X _ {\alpha}, y \rangle h _ {\alpha} = 0. \end{array}
\]

Thus, \(y\) is primitive for \(\rho\) , of weight \(-2\) , which contradicts Prop. 2 of §1, no. 2. This proves (i).

\begin{enumerate}
\def\labelenumi{\alph{enumi})}
\setcounter{enumi}{4}
\tightlist
\item
  Assertion (iii) is now a consequence of c). On the other hand, if \(X_{\alpha}\) is a non-zero element of \(g^{\alpha}\) , the element \(X_{-\alpha}\) constructed in c) is the unique element of \(g^{-\alpha}\) such that \([X_{\alpha}, X_{-\alpha}] = -H_{\alpha}\) since \(\dim g^{-\alpha} = 1\) . The last assertion of (iv) is a consequence of c). Q.E.D.
\end{enumerate}

The notations \(h_{\alpha}\) , \(H_{\alpha}\) , \(s_{\alpha}\) will be retained in what follows. (To define \(h_{\alpha}\) , we take \(\langle\cdot,\cdot\rangle\) equal to the Killing form.) If \(X_{\alpha}\) is a non-zero element of \(g^{\alpha}\) , the isomorphism \(\varphi\) of Th. 1 and the representation \(u \mapsto \operatorname{ad}_{\mathfrak{g}} \varphi(u)\) of \(\mathfrak{sl}(2,k)\) on g will be said to be associated to \(X_{\alpha}\) .

COROLLARY. Let \(\Phi\) be the Killing form of \(\mathfrak{g}\) . For all \(a, b \in \mathfrak{h}\) ,

\[
\Phi (a, b) = \sum_ {\gamma \in \mathrm{R}} \gamma (a) \gamma (b).
\]

Indeed, ad a.ad b leaves each \(g^{\gamma}\) stable, and its restriction to \(g^{\gamma}\) is the homothety with ratio \(\gamma(a)\gamma(b)\) ; if \(\gamma\neq0\) , \(\dim g^{\gamma}=1\) .

PROPOSITION 2. Let \(\alpha, \beta \in \mathbb{R}\) .

\begin{enumerate}
\def\labelenumi{(\roman{enumi})}
\item
  \(\beta(H_{\alpha}) \in \mathbf{Z}\) .
\item
  If \(\Phi\) denotes the Killing form of \(\mathfrak{g}\) , \(\Phi(H_{\alpha}, H_{\beta}) \in \mathbf{Z}\) .
\end{enumerate}

Let \(X_{\alpha}\) be a non-zero element of \(\mathfrak{g}^{\alpha}\) , and let \(\rho\) be the representation of \(\mathfrak{sl}(2,k)\) on \(\mathfrak{g}\) associated to \(X_{\alpha}\) . The eigenvalues of \(\rho(\mathrm{H})\) are 0 and the \(\beta(H_{\alpha})\)

for \(\beta \in \mathbb{R}\) . Hence (i) follows from §1, no. 2, Cor. of Prop. 2. Assertion (ii) follows from (i) and the Cor. of Th. 1. Q.E.D.

Let \(\alpha \in \mathrm{R}\) , \(X_{\alpha}\) a non-zero element of \(\mathfrak{g}^{\alpha}\) , \(X_{-\alpha}\) the element of \(\mathfrak{g}^{-\alpha}\) such that \([X_{\alpha}, X_{-\alpha}] = -H_{\alpha}\) , and \(\rho\) the representation of \(\mathfrak{sl}(2, k)\) on \(\mathfrak{g}\) associated to \(X_{\alpha}\) . Let \(\pi\) be the representation of \(\mathbf{SL}(2, k)\) on \(\mathfrak{g}\) compatible with \(\rho\) ( \(\S 1\) , no. 4, Th. 2). Since \(\operatorname{ad} X_{\alpha}\) is nilpotent (Chap. VII, \(\S 1\) , no. 3, Prop. 10 (iv)), \(\pi(e^{X_+}) = e^{\operatorname{ad} X_{\alpha}}\) is an elementary automorphism of \(\mathfrak{g}\) . Similarly, \(\pi(e^{X_-}) = e^{\operatorname{ad} X_{-\alpha}}\) is an elementary automorphism of \(\mathfrak{g}\) . Hence \(\pi(\mathbf{SL}(2, k)) \subset \operatorname{Aut}_e(\mathfrak{g})\) . We make use of the notation \(\theta(t)\) of \(\S 1\) , no. 5. For \(t \in k^*\) , put

\[
\theta_ {\alpha} (t) = \pi (\theta (t)) = e ^ {\mathrm{ad} t X _ {\alpha}} e ^ {\mathrm{ad} t ^ {- 1} X _ {- \alpha}} e ^ {\mathrm{ad} t X _ {\alpha}}.\tag{1}
\]

Lemma 1. (i) For all \(h \in \mathfrak{h}\) , \(\theta_{\alpha}(t).h = h - \alpha(h)H_{\alpha}\) .

\begin{enumerate}
\def\labelenumi{(\roman{enumi})}
\setcounter{enumi}{1}
\item
  For all \(\beta \in \mathrm{R}\) , \(\theta_{\alpha}(t)(\mathfrak{g}^{\beta}) = \mathfrak{g}^{\beta - \beta (H_{\alpha})\alpha}\) .
\item
  If \(\alpha, \beta \in \mathrm{R}\) , \(\beta - \beta(H_{\alpha})\alpha \in \mathrm{R}\) .
\end{enumerate}

Let \(h \in \mathfrak{h}\) . If \(\alpha(h) = 0\) , \([X_{\alpha}, h] = [X_{-\alpha}, h] = 0\) , so \(\theta_{\alpha}(t).h = h\) . On the other hand, the formulas (5) of §1, no. 5 show that \(\theta_{\alpha}(t).H_{\alpha} = -H_{\alpha}\) . This proves assertion (i). It follows that \(\theta_{\alpha}(t)^2 | \mathfrak{h} = \mathrm{Id}\) . If \(x \in \mathfrak{g}^{\beta}\) and \(h \in \mathfrak{h}\) ,

\[
\begin{array}{r l} [ h, \theta_ {\alpha} (t) x ] & = \theta_ {\alpha} (t). [ \theta_ {\alpha} (t) h, x ] - \beta (\theta_ {\alpha} (t) h). \theta_ {\alpha} (t) x \\ & = (\beta (h) - \alpha (h) \beta (H _ {\alpha})). \theta_ {\alpha} (t) x \\ & = (\beta - \beta (H _ {\alpha}) \alpha) (h). \theta_ {\alpha} (t) x \end{array}
\]

so \(\theta_{\alpha}(t)x\in \mathfrak{g}^{\beta -\beta (H_{\alpha})\alpha}\) . This proves (ii). Assertion (iii) follows from (ii).

THEOREM 2. (i) The set \(\mathrm{R} = \mathrm{R}(\mathfrak{g},\mathfrak{h})\) is a reduced root system in \(\mathfrak{h}^*\) .

\begin{enumerate}
\def\labelenumi{(\roman{enumi})}
\setcounter{enumi}{1}
\tightlist
\item
  Let \(\alpha \in \mathbb{R}\) . The map \(s_{\alpha, H_\alpha}: \lambda \mapsto \lambda - \lambda(H_\alpha)\alpha\) from \(\mathfrak{h}^*\) to \(\mathfrak{h}^*\) is the unique reflection \(s\) of \(\mathfrak{h}^*\) such that \(s(\alpha) = -\alpha\) and \(s(\mathrm{R}) = \mathrm{R}\) . For all \(t \in k^*\) , \(s\) is the transpose of \(\theta_\alpha(t)|\mathfrak{h}\) .
\end{enumerate}

First, R generates \(h^{*}\) , for if \(h \in h\) is such that \(\alpha(h) = 0\) for all \(\alpha \in R\) , then ad h = 0 and hence h = 0 since the centre of g is zero. By definition, \(0 \notin R\) . Let \(\alpha \in R\) . Since \(\alpha(H_{\alpha}) = 2\) , \(s = s_{\alpha,H_{\alpha}}\) is a reflection such that \(s(\alpha) = -\alpha\) . Then \(s(R) = R\) by Lemma 1 (iii), and \(\beta(H_{\alpha}) \in \mathbf{Z}\) for all \(\beta \in R\) (Prop. 2 (i)). This shows that R is a root system in \(h^{*}\) . For all \(h \in h\) and \(\lambda \in h^{*}\) ,

\[
\langle s (\lambda), h \rangle = \langle \lambda - \lambda (H _ {\alpha}) \alpha , h \rangle = \langle \lambda , h - \alpha (h) H _ {\alpha} \rangle = \langle \lambda , \theta_ {\alpha} (t) h \rangle
\]

so \(s\) is the transpose of \(\theta_{\alpha}(t)|\mathfrak{h}\) . Finally, we show that the root system R is reduced. Let \(\alpha \in \mathrm{R}\) and \(y \in \mathfrak{g}^{2\alpha}\) . Since \(3\alpha \notin \mathrm{R}\) (Chap. VI, §1, no. 3, Prop. 8), \([X_{\alpha}, y] = 0\) ; on the other hand, \([X_{-\alpha}, y] \in \mathfrak{g}^{-\alpha + 2\alpha} = \mathfrak{g}^{\alpha} = kX_{\alpha}\) , so \([X_{\alpha}, [X_{-\alpha}, y]] = 0\) ; thus

\[
4 y = 2 \alpha (H _ {\alpha}) y = [ H _ {\alpha}, y ] = - [ [ X _ {\alpha}, X _ {- \alpha} ], y ] = 0
\]

so \(y = 0\) and \(\mathfrak{g}^{2\alpha} = 0\) . In other words, \(2\alpha\) is not a root.

Q.E.D.

Identify \(\mathfrak{h}\) canonically with \(\mathfrak{h}^{**}\) . With the notations of Chap. VI, §1, no. 1, we then have, by Th. 2 (ii),

\[
H _ {\alpha} = \alpha^ {\vee} \quad \text { for   all } \alpha \in \mathrm{R}.\tag{2}
\]

The \(H_{\alpha}\) thus form the root system \(\mathbb{R}^{\vee}\) in \(\mathfrak{h}\) inverse to \(\mathbb{R}\) .

We shall call \(\mathrm{R}(\mathfrak{g},\mathfrak{h})\) the root system of \((\mathfrak{g},\mathfrak{h})\) . The reflections \(s_{\alpha,H_{\alpha}}\) will be denoted simply by \(s_{\alpha}\) . The Weyl group, group of weights, Coxeter number \ldots{} of \(\mathrm{R}(\mathfrak{g},\mathfrak{h})\) are called the Weyl group, group of weights, Coxeter number \ldots{} of \((\mathfrak{g},\mathfrak{h})\) . As in Chap. VI, §1, no. 1, we consider the Weyl group as operating not only on \(h^{*}\) , but also on h by transport of structure, so that \(s_{\alpha} = \theta_{\alpha}(t)|\mathfrak{h}\) . Since the \(\theta_{\alpha}(t)\) are elementary automorphisms of g, we have:

COROLLARY. Every element of the Weyl group of \((\mathfrak{g},\mathfrak{h})\) , operating on h, is the restriction to h of an elementary automorphism of g.

For a converse of this result, see §5, no. 2, Prop. 4.

Remark 1. If \(\mathfrak{h}_{\mathbf{Q}}\) (resp. \(\mathfrak{h}_{\mathbf{Q}}^{*}\) ) denotes the \(\mathbf{Q}\) -vector subspace of \(\mathfrak{h}\) (resp. \(\mathfrak{h}^{*}\) ) generated by the \(H_{\alpha}\) (resp. the \(\alpha\) ), where \(\alpha \in \mathbb{R}\) , then \(\mathfrak{h}\) (resp. \(\mathfrak{h}^{*}\) ) can be identified canonically with \(\mathfrak{h}_{\mathbf{Q}} \otimes_{\mathbf{Q}} k\) (resp. with \(\mathfrak{h}_{\mathbf{Q}}^{*} \otimes_{\mathbf{Q}} k\) ) and \(\mathfrak{h}_{\mathbf{Q}}^{*}\) can be identified with the dual of \(\mathfrak{h}_{\mathbf{Q}}\) (Chap. VI, §1, no. 1, Prop. 1). We call \(\mathfrak{h}_{\mathbf{Q}}\) and \(\mathfrak{h}_{\mathbf{Q}}^{*}\) the canonical \(\mathbf{Q}\) -structures on \(\mathfrak{h}\) and \(\mathfrak{h}^{*}\) (Algebra, Chap. II, §8, no. 1, Def. 1). When we mention \(\mathbf{Q}\) -rationality for a vector subspace of \(\mathfrak{h}\) , for a linear form on \(\mathfrak{h}\) , etc., we shall mean these structures, unless we indicate otherwise. When we mention Weyl chambers, or facets, of R(g, h), we shall work in \(\mathfrak{h}_{\mathbf{Q}} \otimes_{\mathbf{Q}} \mathbf{R}\) or \(\mathfrak{h}_{\mathbf{Q}}^{*} \otimes_{\mathbf{Q}} \mathbf{R}\) , that we shall denote by \(\mathfrak{h}_{\mathbf{R}}\) and \(\mathfrak{h}_{\mathbf{R}}^{*}\) .

Remark 2. The root system \(\mathbb{R}^{\vee}\) in \(\mathfrak{h}\) defines a non-degenerate symmetric bilinear form \(\beta\) on \(\mathfrak{h}\) (Chap. VI, §1, no. 1, Prop. 3), namely the form \((a,b)\mapsto \sum_{\alpha \in \mathbb{R}}\langle \alpha ,a\rangle \langle \alpha ,b\rangle\) . By the Cor. to Th. 1, this form is just the restriction of the Killing form to \(\mathfrak{h}\) . The extension of \(\beta |\mathfrak{h}_{\mathbf{Q}}\times \mathfrak{h}_{\mathbf{Q}}\) to \(\mathfrak{h}_{\mathbf{Q}}\otimes_{\mathbf{Q}}\mathbf{R}\) is positive nondegenerate (Chap. VI, §1, no. 1, Prop. 3). On the other hand, we see that the inverse form on \(\mathfrak{h}^*\) of the restriction to \(\mathfrak{h}\) of the Killing form on \(\mathfrak{g}\) is the canonical bilinear form \(\Phi_{\mathrm{R}}\) of R (Chap. VI, §1, no. 12).

Let \((\mathfrak{g}_{1},\mathfrak{h}_{1})\) , \((\mathfrak{g}_{2},\mathfrak{h}_{2})\) be split semi-simple Lie algebras, \(\varphi\) an isomorphism from \(\mathfrak{g}_{1}\) to \(\mathfrak{g}_{2}\) such that \(\varphi(\mathfrak{h}_{1})=\mathfrak{h}_{2}\) . By transport of structure, the transpose of the map \(\varphi|\mathfrak{h}_{1}\) takes \(R(\mathfrak{g}_{2},\mathfrak{h}_{2})\) to \(R(\mathfrak{g}_{1},\mathfrak{h}_{1})\) .

PROPOSITION 3. Let g be a semi-simple Lie algebra, \(h_{1}\) and \(h_{2}\) splitting Cartan subalgebras of g. There exists an isomorphism from \(h_{1}^{*}\) to \(h_{2}^{*}\) that takes \(\mathrm{R}(\mathfrak{g},\mathfrak{h}_{1})\) to \(\mathrm{R}(\mathfrak{g},\mathfrak{h}_{2})\) .

(For more precise results, see §3, no. 3, Cor. of Prop. 10, and §5, no. 3, Prop. 5).

Let \(k'\) be an algebraic closure of \(k\) , \(\mathfrak{g}' = \mathfrak{g} \otimes_k k'\) , \(\mathfrak{h}_i' = \mathfrak{h}_i \otimes_k k'\) . Then \(\mathrm{R}(\mathfrak{g}', \mathfrak{h}_i')\) is the image of \(\mathrm{R}(\mathfrak{g}, \mathfrak{h}_i)\) under the map \(\lambda \mapsto \lambda \otimes 1\) from \(\mathfrak{h}_i^*\) to \(\mathfrak{h}_i^* \otimes_k k' = \mathfrak{h}_i'^*\) . By Chap. VII, §3, no. 2, Th. 1, there exists an automorphism of \(\mathfrak{g}'\) taking \(\mathfrak{h}_1'\) to \(\mathfrak{h}_2'\) , hence an isomorphism \(\varphi\) from \(\mathfrak{h}_1'^*\) to \(\mathfrak{h}_2'^*\) that takes \(\mathrm{R}(\mathfrak{g}', \mathfrak{h}_1')\) to \(\mathrm{R}(\mathfrak{g}', \mathfrak{h}_2')\) . Then \(\varphi | \mathfrak{h}_1^*\) takes \(\mathrm{R}(\mathfrak{g}, \mathfrak{h}_1)\) to \(\mathrm{R}(\mathfrak{g}, \mathfrak{h}_2)\) , and hence \(\mathfrak{h}_1^*\) to \(\mathfrak{h}_2^*\) . Q.E.D.

In view of Prop. 3, the root system of \((\mathfrak{g},\mathfrak{h})\) depends, up to isomorphism, only on g and not on h. In the same way, the Weyl group, group of weights \ldots{} of \((\mathfrak{g},\mathfrak{h})\) are simply called, by abuse of language, the Weyl group, group of weights \ldots{} of g (cf.~also §5, no. 3, Remark 2). If the Dynkin graph of g is of type \(A_{l}\) , or \(B_{l}\) , \ldots{} (cf.~Chap. VI, §4, no. 2, Th. 3), we say that g is of type \(A_{l}\) , or \(B_{l}\) , \ldots.

Recall that, if \(\alpha\) and \(\beta\) are linearly independent roots, the set of \(j \in \mathbf{Z}\) such that \(\beta + j\alpha \in \mathrm{R}\) is an interval \([-q, p]\) of \(\mathbf{Z}\) containing 0, with \(p - q = -\langle \beta, \alpha^{\vee} \rangle = -\beta(H_{\alpha})\) (Chap. VI, §1, no. 3, Prop. 9).

PROPOSITION 4. Let \(\alpha\) and \(\beta\) be linearly independent roots. Let \(p\) (resp. \(q\) ) be the largest integer \(j\) such that \(\beta + j\alpha\) (resp. \(\beta - j\alpha\) ) is a root.

\begin{enumerate}
\def\labelenumi{(\roman{enumi})}
\item
  The vector subspace \(\sum_{-q\leq j\leq p}\mathfrak{g}^{\beta +j\alpha}\) of \(\mathfrak{g}\) is a simple \(\mathfrak{s}_{\alpha}\) -module of dimension \(p + q + 1\) .
\item
  If \(\alpha + \beta\) is a root, then \([\mathfrak{g}^{\alpha}, \mathfrak{g}^{\beta}] = \mathfrak{g}^{\alpha + \beta}\) .
\end{enumerate}

Let \(X_{\alpha}\) (resp. x) be a non-zero element of \(\mathfrak{g}^{\alpha}\) (resp. \(\mathfrak{g}^{\beta+p\alpha}\) ). Then

\[
\begin{array}{c} {[ X _ {\alpha}, x ] \in \mathfrak {g} ^ {\beta + (p + 1) \alpha} = 0} \\ {[ H _ {\alpha}, x ] = (\beta (H _ {\alpha}) + p \alpha (H _ {\alpha})) x = (- p + q + 2 p) x = (p + q) x.} \end{array}
\]

Thus, \(x\) is primitive of weight \(p + q\) for the representation of \(\mathfrak{sl}(2,k)\) on \(\mathfrak{g}\) associated to \(X_{\alpha}\) ; but the \(\mathfrak{sl}(2,k)\) -module \(\sum_{-q\leq j\leq p}\mathfrak{g}^{\beta +j\alpha}\) is of dimension \(p + q + 1\) ; hence it is simple (§1, no. 2, Prop. 2). If \(\alpha +\beta \in \mathrm{R}\) , then \(p\geq 1\) , so the elements of \(\mathfrak{g}^{\beta}\) are not primitive, and hence \([X_{\alpha},\mathfrak{g}^{\beta}]\neq 0\) . Since \([\mathfrak{g}^{\alpha},\mathfrak{g}^{\beta}]\subset \mathfrak{g}^{\alpha +\beta}\) , we see finally that \([\mathfrak{g}^{\alpha},\mathfrak{g}^{\beta}] = \mathfrak{g}^{\alpha +\beta}\) .

Remark 3. Recall that, by Chap. VI, §1, no. 3, Cor. of Prop. 9, the integer \(p + q + 1\) can only take the values 1, 2, 3, 4.

Remark 4. Let \((\mathfrak{g},\mathfrak{h})\) be a split reductive Lie algebra, \(\mathfrak{c}\) the centre of \(\mathfrak{g},\mathfrak{g}' = \mathcal{D}\mathfrak{g}\) , \(\mathfrak{h}' = \mathfrak{h} \cap \mathfrak{g}'\) . Then \(\mathfrak{h} = \mathfrak{c} \times \mathfrak{h}'\) , and we identify \(\mathfrak{h}'^*\) with a vector subspace of \(\mathfrak{h}^*\) . For any \(\lambda \in \mathfrak{h}^*\) such that \(\lambda \neq 0\) , the primary subspace \(\mathfrak{g}^\lambda\) relative to \(\lambda\) is equal to \(\mathfrak{g}'^{\lambda | \mathfrak{h}'}\) . A non-zero weight of \(\mathfrak{h}\) on \(\mathfrak{g}\) is called a root of \((\mathfrak{g},\mathfrak{h})\) ; every root vanishes on \(\mathfrak{c}\) . Denote by \(\mathrm{R}(\mathfrak{g},\mathfrak{h})\) the set of roots of \((\mathfrak{g},\mathfrak{h})\) ; it can be identified canonically with \(\mathrm{R}(\mathfrak{g}',\mathfrak{h}')\) . Let \(\alpha \in \mathrm{R}(\mathfrak{g},\mathfrak{h})\) . We define \(h_\alpha\) , \(H_\alpha\) , \(\mathfrak{s}_\alpha\) , the isomorphisms \(\mathfrak{sl}(2,k) \to \mathfrak{s}_\alpha\) , and the representations of \(\mathfrak{sl}(2,k)\) on \(\mathfrak{g}\) associated to \(\alpha\) , as in the semi-simple case.

